\emph{Pour chacune des cinq affirmations suivantes, indiquer si elle est \textbf{vraie} ou \textbf{fausse}, en justifiant la réponse.\\
Une réponse non justifiée n’est pas prise en compte.\\
Une absence de réponse n’est pas pénalisée.}

\medskip

\begin{enumerate}
\item On considère la fonction $f$ définie sur l’intervalle $]0~;~ +\infty[$ par :

\[f(x) = \ln (x) - x^2.\]

\textbf{Affirmation 1 :} $\displaystyle \lim_{x \to +\infty} f(x) = -\infty$.

\item On considère l’équation différentielle 

\[(E) : \quad - 2y' + 3y = \sin x + 8\cos x.\]

On considère la fonction $f$ définie sur $\R$ par :

\[f(x) = 2\cos x - \sin x.\]

\textbf{Affirmation 2 :} La fonction $f$ est solution de l’équation différentielle $(E)$.

\item On considère la fonction $g$ définie sur l’intervalle $]0~;~+\infty[$ par :

\[g(x) = \ln (3x + 1) + 8.\]

On considère la suite $(u_n)$ définie par $u_0 = 25$ et pour tout entier naturel $n$ :

\[u_{n+1} = g(u_n).\]

On admet que la suite $(u_n)$ est strictement positive.

\textbf{Affirmation 3 :} La suite $(u_n)$ est décroissante.

\item On considère une fonction affine $h$ définie sur $\R$.

On note $k$ la fonction définie sur $\R$ par $k(x) = x^4 + x^2 + h(x)$.

\textbf{Affirmation 4 :} La fonction $k$ est convexe sur $\R$.

\item Une anagramme d'un mot est le résultat d'une permutation des lettres de ce mot.

Exemple : le mot BAC est possède 6 anagrammes : BAC, BCA, ABC, ACB, CAB, CBA.

\textbf{Affirmation 5 :} Le mot EULER possède 120 anagrammes.
\end{enumerate}
